% Generated draft: requires mathematical and source review.
\documentclass[11pt]{article}
\usepackage[margin=25mm]{geometry}
\usepackage{fontspec}
\setmainfont{Latin Modern Roman}
\usepackage{amsmath,amssymb,mathtools,unicode-math}
\setmathfont{Latin Modern Math}
\usepackage{enumitem,needspace,xurl,hyperref}
\setlength{\parindent}{0pt}
\setlength{\parskip}{5pt}
\setlength{\emergencystretch}{3em}
\newcommand{\N}{\mathbb{N}}
\newcommand{\Z}{\mathbb{Z}}
\newcommand{\Q}{\mathbb{Q}}
\newcommand{\R}{\mathbb{R}}
\newcommand{\C}{\mathbb{C}}
\newcommand{\F}{\mathbb{F}}
\newcommand{\A}{\mathbb{A}}
\newcommand{\PP}{\mathbb{P}}
\newcommand{\E}{\mathbb{E}}
\newcommand{\eps}{\varepsilon}
\newcommand{\dd}{\,\mathrm{d}}
\begin{document}
\section*{UnsolvedMath ID: 701764}
\textbf{Source problem code:} catalogue entry
\subsection{Definitions and mathematical statement}
Use the conventions \(u^{g}=g^{-1}ug\), \([u,v]=u^{-1}v^{-1}uv\), and \([u,v,w]=[[u,v],w]\), with longer commutators also associated to the left. An element \(x\in G\) is a left \(3\)-Engel element if
\[
[g,x,x,x]=1\qquad\text{for every }g\in G.
\]
A \(2\)-group means a group in which every element has finite order a power of \(2\). A group \(G\) is residually finite if, for each \(y\in G\setminus\{1\}\), there are a finite group \(Q\) and a homomorphism \(\varphi:G\to Q\) such that \(\varphi(y)\ne1\). Define the normal closure of \(x\) by
\[
\langle x^{G}\rangle=\langle x^{g}:g\in G\rangle.
\]
A group is locally nilpotent if every finitely generated subgroup is nilpotent. For a group \(H\), write \(\gamma _1(H)=H\) and \(\gamma _{i+1}(H)=[\gamma_i(H),H]\), where the latter denotes the subgroup generated by the indicated commutators. Nilpotency of class at most \(c\) means \(\gamma_{c+1}(H)=\{1\}\).

\textbf{Question (Kourovka Notebook 20.106).} For every residually finite \(2\)-group \(G\) and every \(x\in G\) satisfying
\[
x\ne1,\qquad x^2=1,\qquad [g,x,x,x]=1\quad(g\in G),
\]
is \(\langle x^{G}\rangle\) locally nilpotent? Equivalently, must \(\langle x^{g_1},\ldots,x^{g_m}\rangle\) be nilpotent for every integer \(m\ge1\) and every \(g_1,\ldots,g_m\in G\)? No uniform bound on nilpotency class is requested.

\textbf{Supplied context and comment of 2025.} The source reports that a left \(3\)-Engel element of odd order in any group belongs to the locally nilpotent radical, meaning the largest normal locally nilpotent subgroup, and attributes this result to E. Jabara and G. Traustason (2019).

The comment defines a strong left \(3\)-Engel element \(a\in G\) by the conditions
\[
\gamma_3(\langle a,a^g\rangle)=\{1\},\qquad
\gamma_4(\langle a,a^g,a^h\rangle)=\{1\}
\quad\text{for all }g,h\in G.
\]
It reports that, for elements of odd order, this condition is equivalent to being left \(3\)-Engel. It further reports that, in an arbitrary group, a strong left \(3\)-Engel element \(a\) has locally nilpotent normal closure \(\langle a\rangle^G\), attributing this result to A. Hadjievangelou and G. Traustason (2024). The strong condition is not an additional hypothesis of the question above, and the stated equivalence is qualified by odd order. These contextual results are supplied source assertions, not independently verified conclusions of this draft.
\subsection{Short English statement}
Must the subgroup generated by all conjugates of a left \(3\)-Engel involution in a residually finite \(2\)-group be locally nilpotent? The supplied record labels this question open. Its odd-order results and its 2025 comment about strong left \(3\)-Engel elements are retained as attributed context.
\subsection{Sources}
\begin{itemize}
\item Local UnsolvedMath registry, permanent ID \(701764\), exact source code \texttt{catalogue entry}, title \emph{Kourovka Notebook Problem 20.106}. The supplied record labels the problem \texttt{open} and attributes it to G. Traustason, proposed in 2022. This status is supplied registry information, not an independently verified literature assessment.
\item \emph{The Kourovka Notebook (Unsolved Problems in Group Theory)}, Problem 20.106; the supplied background identifies the 20th issue (2022) and the 21st issue edition (2026), \texttt{arXiv:1401.0300v48}. The supplied collection description names E. I. Khukhro and V. D. Mazurov as editors.
\item E. Jabara and G. Traustason, \emph{Proc. Amer. Math. Soc.} \textbf{147}, no. 5 (2019), 1921--1927. Supplied reference for the assertion about odd-order \(3\)-Engel elements and the locally nilpotent radical.
\item A. Hadjievangelou and G. Traustason, \emph{Proc. Amer. Math. Soc.} \textbf{152}, no. 4 (2024), 1467--1477. Supplied reference in the \emph{Comment of 2025} for local nilpotence of the normal closure of a strong left \(3\)-Engel element.
\end{itemize}
\end{document}
